\subsubsection{Die Isotypie-Zerlegung von N=3 (beide Spuren unabhängig,
identisch)}\label{die-isotypie-zerlegung-von-n3-beide-spuren-unabhuxe4ngig-identisch}

{\def\LTcaptype{none} % do not increment counter
\begin{longtable}[]{@{}lrrrll@{}}
\toprule\noalign{}
Block & Dim & Mult & S-Eigenwert & (C\_S, C\_U) & Rolle \\
\midrule\noalign{}
\endhead
\bottomrule\noalign{}
\endlastfoot
(144, 20) & 2880 & 2 & \textbf{7} & (85, 39) & hell, X-gekoppelt \\
(144, 4̄) & 576 & 2 & \textbf{10} & (85, 15) & hell, X-gekoppelt \\
(16′, 20) & 320 & 2 & \textbf{12} & (45, 39) & hell, X-gekoppelt \\
(16′, 4̄) & 64 & 1 & 0 (ker S) & (45, 15) & χ-Kompositblock \\
(672, 4̄) & 2688 & 1 & - & (165, 15) & dunkel \\
(1200, 20) & 24 000 & 1 & - & (141, 39) & dunkel \\
(560, 20′) & 11 200 & 1 & - & (117, 63) & dunkel \\
\end{longtable}
}

Schur-Funktor-Zerlegung exakt geschält: Sym³(16) = 144⊕\allowbreak{}672, S₂₁(16) =
16′⊕\allowbreak{}144⊕\allowbreak{}1200, Λ³(16) = 560. In Λ³(16,4) multiplizitätsfrei; die drei
hellen Typen kommen je zweimal vor (eine Kopie pro Seite, X-verknüpft).
\textbf{Bemerkenswerter Befund (Parallelsession, von meiner Zerlegung
bestätigt): alle drei dunklen Blöcke tragen denselben Gesamtcasimir 45}
- die eine erhaltene Summeninvariante kann sie prinzipiell nicht
trennen; getrennt gelesen (141/4, 117/4, 165/4 bzw. 39/4, 63/4, 15/4)
trennen sie sofort.

\subsubsection{\texorpdfstring{Das Verfügbarkeits-Theorem
(Parallelsession,
\codeword{symmetry\_availability.py})}{Das Verfügbarkeits-Theorem (Parallelsession, )}}\label{das-verfuxfcgbarkeits-theorem-parallelsession-symmetry_availability.py}

Die gewährte Dynamik (Wörter aus X, N\_b) liegt in der
Verflechteralgebra (Dim Σmᵢ² = \textbf{16}); die 60 Lie-Erzeuger liegen
\textbf{nicht} darin - keine Komposition gewährter Dynamik erzeugt sie.
Die Symmetrie ist ein eigenständiges Primitiv: \emph{gewährt oder nicht
gewährt}. Der gesamte verbleibende symmetrieverträgliche Spielraum ist
\textbf{exakt 2 Dimensionen breit} = die Fähigkeit, die drei dunklen
Blöcke getrennt anzusprechen (= beide Casimire getrennt statt nur ihre
Summe). Die Clock sitzt im Normalisator, nicht im Zentralisator (mit X,
N\_b allein: Algebra 84, Kommutant 240 742 144; über voller Symmetrie
fügt sie nichts hinzu).

\textbf{Antwort auf die §10.3-Frage:} die operative Mehrdeutigkeit ist
auf die binäre Alternative „kontinuierliche Symmetrie verfügbar
ja/nein'' kollabiert - eingeklammert zwischen Kommutant \textbf{7} und
\textbf{1 444 233 216}.

\begin{center}\rule{0.5\linewidth}{0.5pt}\end{center}

\subsection{3. WP-C: Relativistisches Feldwörterbuch (fertig, 2161
exakte Guards,
-OO-stabil)}\label{wp-c-relativistisches-feldwuxf6rterbuch-fertig-2161-exakte-guards--oo-stabil}

Am tatsächlichen Tensor, vor jeder Kontinuumsrechnung:

{\def\LTcaptype{none} % do not increment counter
\begin{longtable}[]{@{}
  >{\raggedright\arraybackslash}p{(\linewidth - 4\tabcolsep) * \real{0.3000}}
  >{\raggedleft\arraybackslash}p{(\linewidth - 4\tabcolsep) * \real{0.4000}}
  >{\raggedright\arraybackslash}p{(\linewidth - 4\tabcolsep) * \real{0.3000}}@{}}
\toprule\noalign{}
\begin{minipage}[b]{\linewidth}\raggedright
Kanal
\end{minipage} & \begin{minipage}[b]{\linewidth}\raggedleft
Norm² pro Zeile
\end{minipage} & \begin{minipage}[b]{\linewidth}\raggedright
Status
\end{minipage} \\
\midrule\noalign{}
\endhead
\bottomrule\noalign{}
\endlastfoot
Skalar, gleichhändig ε\_\{αβ\}ψψ & \textbf{0} (alle 60 Zeilen, exakt) &
\textbf{verboten} \\
Skalar, gepunktet (0,0) & 0 & verboten \\
(1,0) symmetrisch \{I, σˣ, σᶻ\} & je 64, Summe \textbf{192} & erlaubt -
der Vertex ist \textbf{rein (1,0)+h.c.} \\
(0,1) konjugiert & je 64 & erlaubt \\
\end{longtable}
}

Korrektur meines Spec-Fehlers durch den Agenten: σʸ = −iε ist
\emph{anti}symmetrisch; die symmetrische Basis ist \{I, σˣ, σᶻ\}.
Dimensionsidentität C(128,2) = 2080 + 3·2016 = 8128 exakt. Kompositfeld
χ = Vermittler × Weyl: zerfällt (3/2,0)⊕\allowbreak{}(1/2,0), Projektoren Rang 4/2
exakt idempotent, flacher Spin-1/2-Anteil 1/3 (als Konvention
deklariert). Bilineare: 4096 = 1 (Singulett) + 45 + 15 (Adjungierte) +
4035 (Rest, ehrlich unbestimmt). Schur-Folge: kinetischer Operator auf
dem irreduziblen 64er-Multiplett = I₆₄ ⊗ d(p).

\textbf{Konsequenz:} jede relativistische Vervollständigung muss den 60
Vermittlern eine selbstduale Tensor-/feldstärkeartige Lorentzstruktur
((1,0), 3 Komponenten pro internem Label) geben - \textbf{der skalare
Kopplungskanal ist auf dem tatsächlichen Tensor exakt null und darf
nicht wieder verwendet werden}. Die Parallelsession ergänzte eine
minimale skalare Alternative mit extra alternierendem Hilfsindex (128
Komponenten) - ausdrücklich als \emph{zusätzliches} Modell markiert,
nicht als native Lösung.

\begin{center}\rule{0.5\linewidth}{0.5pt}\end{center}

\subsection{4. WP-B: Nativer Grundzustand (Parallelsession v1.6.4 +
meine laufende
Sonde)}\label{wp-b-nativer-grundzustand-parallelsession-v1.6.4-meine-laufende-sonde}

\textbf{Reproduziert (Parallelsession, frische Enumeration statt
Statusübernahme):}

{\def\LTcaptype{none} % do not increment counter
\begin{longtable}[]{@{}rrr@{}}
\toprule\noalign{}
k & Bosonkonfigurationen vollständig enumeriert & ‖Q₊ᵏF‖² \\
\midrule\noalign{}
\endhead
\bottomrule\noalign{}
\endlastfoot
0 & analytisch & 1 \\
1 & 60 Kanäle / 480 Paare & 480 \\
2 & 1 830 & 439 680 \\
3 & 37 820 & 575 078 400 \\
4 & 595 665 & 952 296 652 800 \\
\end{longtable}
}

\textbf{Modellinterner Satz:} für \(0<|g|/\Delta\le 1/20\) ist der
globale Grundzustand Ω \textbf{eindeutig, N=64,
Spin(10)×SU(4)-invariant}; am Prüfpunkt
\(-1{,}158089\,\Delta < E_0 < -1{,}129636\,\Delta\), Lücke
\(> 0{,}007737\,\Delta\). (Der volle Grundzustandsvektor ist damit nicht
ausgerechnet - ehrlich markiert.)

\textbf{Meine Session beigesteuert (Sonde, Teil 1):} die
\textbf{Z4-Zentrumsregel} - Singuletts existieren nur in Sektoren
\(N\equiv 0 \bmod 4\) (SU(4)-Zentrumsphase i\^{}N, niveauunabhängig,
exakt an Basiszuständen geprüft). Das erklärt, warum die Sektorwahl nur
unter N = \ldots, 56, 60, 64, 68, \ldots{} stattfindet; alle anderen
Sektoren haben Casimir-Floors.

\textbf{Noch laufend} (\href{aa1ce128-e016-4c52-ae8f-04d625b04c23}{WP-B
Sonde}): exakte Singulett-Multiplizitäten mult₁ (=1, der R-Zustand) und
mult₂ (Gewicht-0-Unterraum + Casimir-Kern), die exakten
Lanczos-Koeffizienten β₂², β₃² (Erwartung aus der Normsequenz: β₁²=480,
β₂²=439680/480=\textbf{916}, β₃²=299520/229 ≈ 1307,95 - diese
Verhältnisform gilt genau dann, wenn die Singuletts pro Niveau
eindimensional bleiben, also ist mult₂ der Zwickel), Casimir-Floor
c\_min=56 auf Λ², Sektorfloors N=56\ldots72, Gruppen-Census. Damit prüfe
ich die Kernzahlen der Parallelsession unabhängig nach.

\begin{center}\rule{0.5\linewidth}{0.5pt}\end{center}

\subsection{5. WP-D: Geladene Antwort auf dem Grundzustand
(Parallelsession
v1.6.4)}\label{wp-d-geladene-antwort-auf-dem-grundzustand-parallelsession-v1.6.4}

Genau die geforderte „dieselbe geladene Antwort auf dem abgesicherten
Zustand'':

\begin{itemize}
\tightlist
\item
  \textbf{Bewegungsgleichung (Operatoridentität, endlicher
  Teilchenkern):} \([H,f_r]=-gD_r\),
  \(D_r=\sum_{A,j}(M_A)_{rj}b_Af_j^\dagger\), Ladung −1 (nicht das
  frühere χ† mit Ladung +3). Kompositnorm:
  \(\langle\{D_r,D_s^\dagger\}\rangle=\delta_{rs}S\),
  \(S=15-\tfrac{7}{32}\bar b\).
\item
  \textbf{Exakte Spektralmomente:} \(Z_+=\bar b/32\),
  \(Z_-=1-\bar b/32\), \(a=(\Delta\bar b-E_0)/64\);
  \(m_0=1,\ m_1=0,\ m_2=g^2S,\ m_3=g^2(\Delta S+7a)\) (m₃ ist der eigene
  neue Schritt; unabhängig an der exakt lösbaren N=4/N=5-Antwort
  kontrolliert).
\item
  \textbf{Verschärfte rationale Schranken} (Polynom \(P(b)>0\),
  Wurzelisolation): \(0{,}842846<\bar b<1{,}245656\).
\item
  \textbf{Polsatz (v1.6.3, reproduziert):} isolierte Entnahmelinie,
  64-fach entartet (duale 64), \(Z_{\text{low}}>0{,}880076280689\ldots\)
  - \textbf{mehr als 88 \% des Spektralgewichts pro Mode} -, Polfenster
  \(0{,}007737<\epsilon/\Delta<0{,}039079764\), übrige Entnahmeenergien
  \(>0{,}379636\,\Delta\), Additionen \(>0{,}329636\,\Delta\).
\item
  \textbf{Zwei-Linien-Ausschluss:} eine exakte Zwei-Linien-Antwort
  erzwingt \(\epsilon_+-\epsilon_-=\Delta\), aber
  \(m_3/m_2=\Delta+7a/S>\Delta\) - Widerspruch. \(\Sigma_2\not\equiv0\)
  ist bewiesen; die Reststruktur ist unvermeidlich.
\end{itemize}

\begin{center}\rule{0.5\linewidth}{0.5pt}\end{center}

\subsection{6. Kreuzvalidierung der beiden
Spuren}\label{kreuzvalidierung-der-beiden-spuren}

{\def\LTcaptype{none} % do not increment counter
\begin{longtable}[]{@{}
  >{\raggedright\arraybackslash}p{(\linewidth - 6\tabcolsep) * \real{0.2500}}
  >{\raggedright\arraybackslash}p{(\linewidth - 6\tabcolsep) * \real{0.2500}}
  >{\raggedright\arraybackslash}p{(\linewidth - 6\tabcolsep) * \real{0.2500}}
  >{\raggedright\arraybackslash}p{(\linewidth - 6\tabcolsep) * \real{0.2500}}@{}}
\toprule\noalign{}
\begin{minipage}[b]{\linewidth}\raggedright
Befund
\end{minipage} & \begin{minipage}[b]{\linewidth}\raggedright
Meine Session
\end{minipage} & \begin{minipage}[b]{\linewidth}\raggedright
Parallelsession
\end{minipage} & \begin{minipage}[b]{\linewidth}\raggedright
Status
\end{minipage} \\
\midrule\noalign{}
\endhead
\bottomrule\noalign{}
\endlastfoot
N=3-Zerlegung (7 Blöcke, Mults 2,2,2,1,1,1,1) & ✓ & ✓ &
\textbf{identisch} \\
Casimirwerte (je Block) & (45,15), (85,15), \ldots{} & C = (C\_S+C\_U)/4
= 15, 25, 45, 21, 31, 45, 45 & \textbf{identisch} (zwei Normierungen) \\
Kommutant S1 (volle Symmetrie) & 7 & 7 & \textbf{identisch} \\
Kommutant S0 (X, N\_b) & \st{1 452 961 792} → \textbf{1 444 233 216} & 1
444 233 216 & \textbf{Diskrepanz gefunden + von mir korrigiert} \\
Skalarkanal exakt null & ✓ (60 Zeilen, beide Chiralitäten) & ✓ &
\textbf{identisch} \\
Grundzustandsnormen & Sonde läuft (β₁²=480 ✓) & 1, 480, 439 680, 575 078
400, 952 296 652 800 & Abgleich läuft \\
\end{longtable}
}

\begin{center}\rule{0.5\linewidth}{0.5pt}\end{center}

\subsection{7. Offen / nächste Schritte (in der Reihenfolge der
Direktive)}\label{offen-nuxe4chste-schritte-in-der-reihenfolge-der-direktive}

\begin{enumerate}
\def\labelenumi{\arabic{enumi}.}
\tightlist
\item
  \textbf{Sonde abschließen} (läuft): mult₂-Zertifikat, β-Kette,
  Sektorfloors → unabhängige Reproduktion der v1.6.4-Kernzahlen.
\item
  \textbf{Volle 33-Niveau-Kette:} falls mult\_k=1 weiter trägt, ist die
  Singulett-Kette durch die Normsequenz ‖Q₊ᵏF‖² (k≤32) bestimmt - dafür
  ist die Symmetriespur-Resummation nötig (die „vier Symmetriespuren''
  des Workers), nicht rohe Enumeration.
\item
  \textbf{Σ₂ kontrollieren} (erster Norm-/Momentkoeffizient, N=63/64/65
  gemeinsam) - die v1.6.4-Summenregeln sind die Akzeptanztests.
\item
  \textbf{Feldtyp an der Quelle entscheiden:} (1,0)-Vermittler mit
  Kinetik \emph{oder} nachgewiesene zweite Komponente für den skalaren
  Typ.
\item
  \textbf{Präparationslücke:} zahlenerhaltende Kontrollen erreichen N=64
  aus dem leeren Zustand nicht (Parallelsession, §8.3);
  imaginärzeitliches \(e^{-\tau H}F\) konvergiert mathematisch, braucht
  aber Instrumentenherleitung.
\item
  \textbf{Räumliche Skalierung: bewusst NICHT begonnen} - gemäß deiner
  Direktive trägt sie erst nach Zustand + Feldwörterbuch.
\item
  \textbf{Abschluss meines Contracts:} \codeword{replay.py} +
  \codeword{RESULTS.md} + README, wenn die Sonde zurück ist (inkl. des
  Erratums zum v1.6.2-Pluszeichen, das die Parallelsession dokumentiert
  hat).
\end{enumerate}

Kein T1-T8-Tor geschlossen, keine Promotion nach
\codeword{verification/}, kein Ledger/Paper-Eingriff - alles bleibt im
\codeword{experiments/}-Firewall.

\begin{center}\rule{0.5\linewidth}{0.5pt}\end{center}

